\section{Curvature blocks and their covariant derivatives}
\label{sec:preliminaries}

Our first task is to express sectional positivity and the differential
Bianchi identity in a common set of variables. The curvature decomposition
is classical; see Derdzi\'nski~\cite[Section~2]{Derdzinski1983}.
We give the algebra and fix its normalizations, since the later Hessian
estimate depends on the precise representation constants.
Rescale the metric so that $\operatorname{Ric}=3g$, and write
\[
 d\nu=\frac{d\operatorname{vol}_g}{\operatorname{Vol}(M)},
 \qquad \mathbf E f=\int_M f\,d\nu.
\]
The Laplacian $\Delta$ is the trace of the second covariant derivative.
Our curvature convention is
\[
 [\nabla_i,\nabla_j]\xi_k=R_{ijka}\xi_a
\]
for a covector, with $R_{ijij}>0$ on positively curved orthonormal
two-planes. All tensor connections are induced by the Levi--Civita
connection. We take covariant derivatives before choosing any
pointwise eigenbasis.

\subsection{The two Weyl blocks}

The Hodge-star eigenspaces $\Lambda^+$ and $\Lambda^-$ have eigenvalues
$+1$ and $-1$, respectively. On an oriented Euclidean four-space, choose
triples of skew orthogonal
complex structures $I_1,I_2,I_3$ and $J_1,J_2,J_3$ satisfying
\[
 I_1I_2=I_3=-I_2I_1,\qquad
 J_1J_2=J_3=-J_2J_1,
\]
and their cyclic analogues, with every $I_s$ commuting with every $J_t$.
The two triples can be constructed from left and right quaternionic
multiplication. As two-forms, $I_s/\sqrt2$ and $J_s/\sqrt2$ are
orthonormal bases of $\Lambda^+$ and $\Lambda^-$, respectively.
Vector and covector indices are identified by the metric.

\begin{lemma}[Curvature blocks and sectional positivity]
\label{lem:curvature-blocks}
There are symmetric trace-free endomorphisms
$T\in\operatorname{Sym}_0(\Lambda^+)$ and
$U\in\operatorname{Sym}_0(\Lambda^-)$ such that the curvature operator is
\[
 \mathcal R_{\Lambda^2}
 =\begin{pmatrix}\operatorname{Id}+T&0\\
                  0&\operatorname{Id}+U\end{pmatrix}.
\]
At a point where these endomorphisms are diagonal, with eigenvalues
$t_s,u_s$, the curvature tensor is
\[
 R_{ikjl}
 =\frac12\sum_s(1+t_s)I_{s,ik}I_{s,jl}
  +\frac12\sum_s(1+u_s)J_{s,ik}J_{s,jl}.
\]
Put
\[
 x=-\min_s t_s,\qquad y=-\min_s u_s,\qquad
 v=\frac{|T|}{\sqrt6},\qquad b=\frac{|U|}{\sqrt6},
\]
where the matrix norms are Hilbert--Schmidt norms. Then
\[
 \sec_g>0\quad\Longleftrightarrow\quad x+y<2,
 \qquad
 v\le x\le2v,\qquad b\le y\le2b.
\]
Both Weyl blocks, viewed as curvature four-tensors, have zero
divergence. In particular, for $Z=\nabla T$, whose matrix entries are
one-forms, the differential Bianchi constraint is
\[
 \sum_s I_s Z_{sr}=0\qquad(r=1,2,3).
\]
\end{lemma}

\begin{proof}
The first Bianchi identity says that the traces of the two diagonal
blocks of the curvature operator agree. Their sum is half the scalar
curvature, hence is $6$. Contraction identifies the off-diagonal block
with the trace-free Ricci tensor: its nine entries give the coefficients
in the independent symmetric trace-free matrices $I_sJ_t$.
Thus the Einstein condition makes this block zero, and each diagonal
block has trace $3$. Expanding in the orthonormal two-form bases gives
the displayed curvature tensor.

A unit simple two-form has self-dual and anti-self-dual parts of norm
$1/\sqrt2$. Conversely, a unit two-form whose two parts have these
norms has exterior square zero, and is simple: in its orthogonal
normal form $\alpha=a\,e^1\wedge e^2+b\,e^3\wedge e^4$, the condition
$\alpha\wedge\alpha=0$ is $ab=0$. Therefore all pairs of plus and minus
parts of norm $1/\sqrt2$ occur. The smallest sectional curvature is
$1-(x+y)/2$, proving the equivalence.

For the norm comparison, write the eigenvalues of $T$ as
$-x,a,x-a$. Minimality of $-x$ gives $-x\le a\le2x$, and hence
\[
 \frac32x^2
 \le x^2+a^2+(x-a)^2
 \le6x^2.
\]
Division by $6$ gives $v\le x\le2v$. The proof for $U$ is identical.

Contracting the differential Bianchi identity shows that the curvature
tensor is divergence-free when Ricci is parallel. Project its last
pair of indices onto $\Lambda^+$ or $\Lambda^-$; these projections
are parallel. Subtracting the parallel scalar-curvature contribution
proves divergence freedom of each Weyl block. In a quaternionic basis,
the resulting constraint on $\nabla T$ is exactly the stated identity.
\end{proof}

\subsection{Algebra of the derivative representations}

Let $\Sigma_s$ denote the infinitesimal action of $I_s$ on the plus
indices of a tensor representation. On $\Lambda^+$ it is rotation
about axis $s$ with speed $2$. A complexified irreducible representation
has \emph{spin $j$} when
\[
 \sum_{s=1}^3\Sigma_s^2=-4j(j+1)\operatorname{Id}.
\]
Pairs $(j,\ell)$ specify the spins of the two commuting rotation
algebras. This notation describes associated tensor bundles of the
oriented orthonormal frame bundle; it requires no spin structure.

\begin{lemma}[Representation rules]
\label{lem:spin-algebra}
A spin-$j$ representation has weights
$-j,-j+1,\ldots,j$ for $\Sigma_3/(2i)$, each once, and
\[
 H_j\otimes H_\ell
 \cong H_{j+\ell}\oplus H_{j+\ell-1}\oplus\cdots\oplus H_{|j-\ell|}.
\]
In particular, the tangent, plus two-form, and plus trace-free
symmetric-matrix representations have types
\[
 \left(\frac12,\frac12\right),\qquad (1,0),\qquad (2,0).
\]
The Casimir operators and the orthogonal projections onto these
irreducible types are parallel bundle endomorphisms.
\end{lemma}

\begin{proof}
The generators are skew-Hermitian and satisfy
$[\Sigma_1,\Sigma_2]=2\Sigma_3$, cyclically. Diagonalize
$\Sigma_3/(2i)$ and take the usual conjugate raising and lowering
operators, which shift its weight by one. If a unit highest vector
has weight $j$, commuting the raising and lowering operators gives
the squared lowering coefficient at weight $j-r$ as
\[
 (r+1)(2j-r).
\]
Nonnegativity and termination of the chain imply that $2j$ is a
nonnegative integer and give the listed weights and Casimir.
Taking invariant orthogonal complements decomposes an arbitrary
representation. Multiplying the weight lists, and successively
removing the chain of highest remaining weight, proves the tensor
product rule with multiplicity one. Commuting actions operate on the
corresponding multiplicity spaces. The same chain description shows
that an intertwiner between single irreducible copies is a scalar.

The actions on the tangent space and on $\Lambda^+$ give the first
two types. The symmetric square of the spin-one representation is
spin two plus the scalar trace, giving the third. Finally, these
splittings are defined by invariant algebraic operators. The
connection induced by the orthonormal frame bundle preserves every
such operator. Thus the projections are parallel even though a
chosen local frame, or a curvature eigenbasis, need not be parallel.
\end{proof}

Write $F=\operatorname{Sym}_0(\Lambda^+)$. For a symmetric trace-free
matrix $A$, define the self-adjoint operator on $F$
\[
 B_F(A)S=\frac{AS+SA}{2}
         -\frac{\operatorname{tr}(AS)}3\operatorname{Id}.
\]
It is characterized by
$\langle B_F(A)S,S\rangle=\operatorname{tr}(AS^2)$.
On any plus spin-$j$ representation with $j\ge1$, define
\[
 Q_j(A)=\frac{\sum_{r,s}A_{rs}\Sigma_r\Sigma_s}{4j(2j-1)}.
\]
For diagonal $A$ the numerator is $\sum_s a_s\Sigma_s^2$.

\begin{lemma}[Curvature-weighted operators]
\label{lem:quadratic-curvature-action}
If $A$ has smallest eigenvalue $-x$, then
\[
 Q_j(A)\le \frac{x(j+1)}{2j-1}\operatorname{Id}.
\]
Moreover, $Q_2(A)=B_F(A)$. The compression of $B_F(A)$, acting on an
$F$ factor, to a highest-spin coupling of total plus spin $j$ is
$Q_j(A)$, provided the operator acts trivially on any additional
irreducible minus factor.
\end{lemma}

\begin{proof}
In a diagonal basis, $a_s+x\ge0$ and every $\Sigma_s^2$ is
nonpositive. Thus
\[
 \sum_s a_s\Sigma_s^2
 \le -x\sum_s\Sigma_s^2
 =4xj(j+1)\operatorname{Id},
\]
which proves the bound.

To identify $Q_2$, diagonalize $S=\operatorname{diag}(d_1,d_2,d_3)$.
The speed-two matrix action gives
$|\Sigma_sS|^2=8(d_j-d_k)^2$, with $s,j,k$ cyclic; these three
derivatives are mutually orthogonal. Consequently, allowing $A$
to be nondiagonal in this basis,
\[
 \begin{aligned}
 \langle Q_2(A)S,S\rangle
 &=-\frac13\sum_s A_{ss}(d_j-d_k)^2\\
 &=-\frac13\sum_s A_{ss}
       \bigl(2\operatorname{tr}(S^2)-3d_s^2\bigr)
 =\operatorname{tr}(AS^2).
 \end{aligned}
\]
Here trace-freeness of both matrices gives the middle identity and
removes the scalar term. Polarization proves the operator equality.

By Lemma~\ref{lem:spin-algebra}, spin two occurs once in
$\operatorname{End}(H_j)\cong H_j\otimes H_j$. Hence an equivariant
operator family linear in $A$ is unique up to scalar. The same is
true with an irreducible minus factor when the operators commute
with its action. Test the scalar at
$A=\operatorname{diag}(-1,-1,2)$ on a highest vector. At weight $m$,
\[
 Q_j(A)=\frac{j(j+1)-3m^2}{j(2j-1)},
\]
so its highest-weight eigenvalue is $-1$. The $F$ highest vector
also has $B_F(A)$ eigenvalue $-1$. A product of highest vectors is
the highest vector of the stretched coupling, proving the claimed
compression. For the same reason, compression of a single generator
from a spin-$\ell$ factor to a stretched spin-$j$ product is
$\ell/j$ times the total generator.
\end{proof}

\begin{proposition}[Differential equations for the Weyl blocks]
\label{prop:weyl-derivative-equations}
The tensor $Z=\nabla T$ belongs to the parallel subbundle
$G$ of type $(5/2,1/2)$. On this bundle, put
$B_{\mathrm{op}}=Q_{5/2}(T)$. Then
\[
 \Delta T=6(T-N),\qquad
 N=(T^2)_0=T^2-2v^2\operatorname{Id},
\]
and
\begin{equation}
 \Delta Z=(13-20B_{\mathrm{op}})Z.
 \tag{1}\label{eq:1}
\end{equation}
Every statement has an exchanged counterpart for $U$.
\end{proposition}

\begin{proof}
The derivative space $T^*M\otimes F$ has types
$(5/2,1/2)$ and $(3/2,1/2)$. The divergence takes values in a space
whose plus spins are at most $3/2$. It is nonzero: a nonzero covector
tensored with $\operatorname{diag}(1,-1,0)$ has nonzero divergence.
It therefore vanishes on the first summand and is injective on the
second, by irreducibility. Lemma~\ref{lem:curvature-blocks} puts $Z$
in the first summand.

Let
\[
 C=\sum_s I_s\otimes\Sigma_s
\]
on $T^*M\otimes F$. Expansion of the combined Casimir gives
$C=-4$ on $G$ and $C=6$ on its complement. In particular
$P_G=(6-C)/10$ is a parallel orthogonal projector. The bundle
curvature on $F$ is
\[
 F_{ik}=\frac12\sum_s(1+t_s)I_{s,ik}\Sigma_s
\]
at a point of diagonalization. Taking the covariant divergence of
$CZ=-4Z$, using skewness of $I_s$ to take a commutator and
$\sum_{ik}I_{s,ik}I_{r,ik}=4\delta_{sr}$, gives
\[
 -4\Delta T=\sum_s(1+t_s)\Sigma_s^2T
            =-24T+24B_F(T)T=-24(T-N).
\]

To differentiate this equation, note first that
$\nabla N=2B_F(T)Z$. For each derivative index $j$, the covariant
commutation formula reads
\[
 \Delta Z_j
 =\nabla_j\Delta T+\nabla^i(F_{ij}T)
                    +F_{ij}\nabla^iT
                    +\operatorname{Ric}_{ja}Z_a.
\]
The divergence of $F$ vanishes because it is the induced action of
the divergence-free Riemann curvature. Thus the differentiated
equation contributes $6Z-12B_F(T)Z$, the Ricci term contributes
$3Z$, and the remaining bundle term is
$-\sum_s(1+t_s)I_s\Sigma_sZ$. Its constant part is $-CZ=4Z$.

The left side already lies in $G$, since $P_G$ is parallel and
$P_GZ=Z$. We may therefore compress the right side to $G$,
although its individual summands need not preserve $G$.
Lemma~\ref{lem:quadratic-curvature-action} gives
$P_GB_F(T)P_G=Q_{5/2}(T)$. On $G$, the total plus generators are
$I_s+\Sigma_s$, so
\[
 40Q_{5/2}(T)
 =P_G\sum_s t_s(I_s+\Sigma_s)^2P_G.
\]
The $I_s^2$ term vanishes by $\sum_s t_s=0$, and the $\Sigma_s^2$
term compresses to $24Q_{5/2}(T)$. Hence
\[
 P_G\sum_s t_s I_s\Sigma_sP_G=8Q_{5/2}(T).
\]
The scalar contributions sum to $13$, while the two curvature
contributions sum to $-20Q_{5/2}(T)$, proving \eqref{eq:1}.
For a nondiagonal $T$ the same calculation uses the symmetric
contraction $\sum_{r,s}T_{rs}(I_r+\Sigma_r)(I_s+\Sigma_s)$:
its $I_rI_s$ part is zero by symmetry and trace-freeness.
Thus the computation is tensorial and entails no differentiation
of a moving eigenframe.
\end{proof}

\subsection{The gradient parametrization and the norm equation}

Write
\[
 p=\nabla v,\qquad q=\nabla b,\qquad
 e=\frac{|\nabla T|^2}{6},\qquad
 \bar e=\frac{|\nabla U|^2}{6}.
\]
The norms $v,b$ are Lipschitz. Their gradients are understood almost
everywhere and vanish almost everywhere on the respective zero sets.

Where $v>0$, diagonalize $T$ and put, cyclically in $i,j,k$,
\[
 t_i=vs_i,\qquad
 \widetilde s_i=\frac{s_j-s_k}{\sqrt3},\qquad
 \sigma=\frac16\sum_i s_i^3,\qquad
 \rho=\frac16\sum_i s_i^2\widetilde s_i.
\]
Here $\sum s_i=0$ and $\sum s_i^2=6$. On the plane of triples with
zero sum, the map $s\mapsto\widetilde s$ is a skew-adjoint orthogonal
quarter-turn. Thus
\[
 \sum_i\widetilde s_i=0,\qquad
 \sum_i\widetilde s_i^2=6,\qquad
 \sum_i s_i\widetilde s_i=0,\qquad
 s_i^2+\widetilde s_i^2=4.
\]
For later use, these identities imply
\[
 \sigma^2+\rho^2=1,\qquad
 \frac{N}{v^2}
   =\operatorname{diag}(\sigma s_i+\rho\widetilde s_i),
 \qquad
 \frac16\sum_i\widetilde s_i^3=-\rho.
\]
Indeed, the trace-free triple $n_i=s_i^2-2$ has squared norm $6$:
use $\sum_i s_i^4=(\sum_i s_i^2)^2/2=18$.
Its coefficients in the orthogonal basis $s,\widetilde s$ are
$\sigma,\rho$, proving the first two assertions. The last follows
by factoring the alternating cubic
$\sum_i s_i^2(s_j-s_k)$. If $r=x/v$, the other two eigenvalues
besides $-r$ have sum $r$ and product $r^2-3$, so
\[
 1\le r\le2,\qquad
 \sigma=\frac{3r-r^3}{2}.
\]
Normalized quantities are used only where the corresponding block
is nonzero. Polynomial products such as
$\sigma v^3=\operatorname{tr}(T^3)/6$ have their natural extensions
across its zero set. Bars denote the analogous quantities for $U$.

The next formula gives, in particular, the refined Kato inequality
$|\nabla T|^2\ge(5/3)|\nabla|T||^2$ for a half-Weyl block.
Gursky--LeBrun derive this estimate from the same irreducible Bianchi
gradient constraint~\cite[Lemma~4, equations~(10)--(12)]{GurskyLeBrun1999};
Calderbank--Gauduchon--Herzlich give its general representation-theoretic
setting~\cite[equation~(7.1)]{CalderbankGauduchonHerzlich2000}.
We retain all three vector parameters: the excess over the norm-gradient
energy will be needed in the coupled estimate.

\begin{lemma}[Parametrization of the Bianchi gradient space]
\label{lem:gradient-parametrization}
There are unique vectors $k,l$ such that
\begin{equation}
 \begin{gathered}
 Z_{ii}=s_i p-\sqrt3\,\widetilde s_i k,\qquad
 Z_{jk}=-I_i\left(l+s_i k+\frac{\widetilde s_i}{\sqrt3}p\right),
 \qquad (i,j,k)\ \text{cyclic},\\
 e=\frac53|p|^2+5|k|^2+|l|^2.
 \end{gathered}
 \tag{2}\label{eq:2}
\end{equation}
More generally, for any fixed normalized trace-free diagonal triple
$s$, the same formulas parametrize the entire algebraic space $G$
by three arbitrary real four-vectors.
\end{lemma}

\begin{proof}
Write $D_i=Z_{ii}$ and $a_i=I_iZ_{jk}$. Quaternion multiplication
turns the three divergence equations into
$D_i+a_j-a_k=0$. Trace-freeness and the orthogonal basis
$s,\widetilde s$ give uniquely
\[
 p=\frac16\sum_i s_iD_i,\qquad
 k=-\frac1{6\sqrt3}\sum_i\widetilde s_iD_i,\qquad
 D_i=s_ip-\sqrt3\,\widetilde s_i k.
\]
Since
$s_j-s_k=\sqrt3\,\widetilde s_i$ and
$\widetilde s_j-\widetilde s_k=-\sqrt3\,s_i$,
the vectors $s_i k+\widetilde s_i p/\sqrt3$ have exactly the
differences required of the $a_i$. Their difference from $a_i$ is
therefore a common vector $l$, proving the formulas and their
converse. This also establishes their full algebraic scope.

Orthogonality of $s,\widetilde s$ gives
\[
 \sum_i|Z_{ii}|^2=6|p|^2+18|k|^2,\qquad
 2\sum_{j<k}|Z_{jk}|^2=4|p|^2+12|k|^2+6|l|^2.
\]
The factor two in the second sum is the Hilbert--Schmidt multiplicity
of off-diagonal entries. Division by six proves the norm identity.
For $Z=\nabla T$,
$\nabla|T|^2=2v\sum_i s_iZ_{ii}=12vp$; comparison with
$|T|^2=6v^2$ identifies the parameter $p$ with $\nabla v$.
\end{proof}

The pointwise norm equation alone is insufficient for integration:
a Weyl block may vanish. We next prove that its regularized Laplacians
converge to an actual integrable density, so no measure supported on the
zero set is lost.

\begin{proposition}[Weak norm equation, including the zero set]
\label{prop:weak-norm-equation}
The function $e/v$, defined to be zero on $\{v=0\}$, belongs to
$L^1(M)$, and
\begin{equation}
 \Delta v=6v(1-\sigma v)+\frac{e-|p|^2}{v}
 \tag{3}\label{eq:3}
\end{equation}
holds as an equality of distributions with an $L^1$ density.
The quotient is zero almost everywhere on the zero set.
This identity may be tested against bounded Lipschitz functions.
The corresponding assertions hold for $b$.
\end{proposition}

\begin{proof}
On $\{v>0\}$, pair the equation for $\Delta T$ with $T$ to obtain
\[
 \Delta(v^2)=12(v^2-\sigma v^3)+2e.
\]
This is globally an identity of smooth functions. Put
$v_\varepsilon=\sqrt{v^2+\varepsilon^2}$. Direct differentiation gives
\[
 \Delta v_\varepsilon
 =\frac{6(v^2-\sigma v^3)}{v_\varepsilon}
   +q_\varepsilon,\qquad
 q_\varepsilon=\frac e{v_\varepsilon}
             -\frac{v^2|p|^2}{v_\varepsilon^3}.
\]
The refined derivative estimate in \eqref{eq:2} yields
\[
 \frac25\frac e{v_\varepsilon}
 \le q_\varepsilon\le\frac e{v_\varepsilon}.
\]
These assertions hold almost everywhere, including the zero set:
every smooth component of $T$ has derivative zero almost everywhere
on its zero set, hence $Z=0$ almost everywhere on $\{T=0\}$.
One can see this elementary fact at density points of a level set;
a nonzero derivative would instead make that scalar level set a
local hypersurface.

The first term in $\Delta v_\varepsilon$ is uniformly bounded in
absolute value by $6(v+v^2)$. Its integral is the negative of the
integral of $q_\varepsilon$. Consequently
$\mathbf E(e/v_\varepsilon)$ is uniformly bounded.
Monotone convergence proves the asserted $L^1$ bound for $e/v$.
Off the zero set, $q_\varepsilon$ is dominated by $e/v$ and tends
to $(e-|p|^2)/v$; on the zero set both are zero almost everywhere.
Dominated convergence therefore gives an $L^1$ limit, rather than
only a weak limit of measures. The zeroth-order term converges in
$L^1$ as well. Since $v_\varepsilon\to v$ uniformly, their
Laplacians converge distributionally, proving \eqref{eq:3} without
a residual singular measure.

Finally, approximate a bounded Lipschitz test function uniformly and
in $H^1$ by smooth functions, using local mollification and a
partition of unity. Uniform convergence controls its pairing with
the $L^1$ right side, and $H^1$ convergence controls its gradient
pairing with $\nabla v$. This proves the asserted testing rule.
\end{proof}
